% ch7_d (书页 311-322 / p326-p337)
%--C-TAIL--

% ==== tail：7.6 节收尾（承接书页 310，从 (7.168) 起），置于下一个 \section 之前 ====

% ---- 书页 311 / p326 ----

把 $\cos^2$ 项在几个波长上取平均，得到
\begin{equation*}
\left\langle \cos^2(k_\lambda x^\lambda) \right\rangle = \frac{1}{2}.
\tag{7.168}
\end{equation*}
为简单起见，我们可以取波沿 $z$ 轴传播，于是
\begin{equation*}
k_\lambda = \left(-\omega,\, 0,\, 0,\, \omega\right)
\tag{7.169}
\end{equation*}
负号来自把 $k^\lambda$ 的一个指标降下来；再利用 (7.109)，得到
\begin{equation*}
C_{\rho\sigma}C^{\rho\sigma} = 2(h_+^2 + h_\times^2).
\tag{7.170}
\end{equation*}
在引力波文献中，更常见的做法是用普通频率 $f = \omega/2\pi$ 而不是角频率 $\omega$ 来表达可观测量。把上面的结果合在一起，就得到
\begin{equation*}
t_{\mu\nu} = \frac{\pi}{8G}f^2(h_+^2 + h_\times^2)
\begin{pmatrix}
1 & 0 & 0 & -1 \\
0 & 0 & 0 & 0 \\
0 & 0 & 0 & 0 \\
-1 & 0 & 0 & 1
\end{pmatrix}.
\tag{7.171}
\end{equation*}
正如我们在下一节将要讨论的，我们有望在地球上观测到的典型引力波源，频率介于 $10^{-4}$ 与 $10^4\,\mathrm{Hz}$ 之间，振幅 $h \sim 10^{-22}$。因此，把沿 $z$ 方向的能量流 $-T_{0z}$ 用数量级的形式表示出来是有用的：
\begin{equation*}
-T_{0z} \sim 10^{-4}\left(\frac{f}{\mathrm{Hz}}\right)^{2}\frac{\left(h_+^2 + h_\times^2\right)}{\left(10^{-21}\right)^{2}}\ \frac{\mathrm{erg}}{\mathrm{cm}^2\cdot\mathrm{s}}.
\tag{7.172}
\end{equation*}
这就是原则上每秒钟可以沉积到探测器每平方厘米中的能量。正如 Thorne\footnote{K.S.~Thorne, 载于 \emph{Three Hundred Years of Gravitation}, Cambridge: Cambridge University Press, 1987。}所指出的，这实际上是相当可观的能量流，尤其是在频率范围的高端。作为比较，处于宇宙学距离上的超新星，其峰值电磁流量约为 $10^{-9}\,\mathrm{erg/cm^2/s}$；然而引力波信号只持续几毫秒，而可见的电磁信号却能延续数月。

现在让我们用引力波能量-动量张量的公式，来计算一个按四极公式 (7.140) 发出引力辐射的系统的能量损失率。在一个常值时间的曲面 $\Sigma$ 上所包含的引力辐射总能量定义为
\begin{equation*}
E = \int_{\Sigma} t_{00}\, d^3x,
\tag{7.173}
\end{equation*}
而辐射到无穷远处的总能量则可以表示为

% ---- 书页 312 / p327 ----

\begin{equation*}
\Delta E = \int P\, dt,
\tag{7.174}
\end{equation*}
其中功率 $P$ 为
\begin{equation*}
P = \int_{S^2_\infty} t_{0\mu}n^\mu r^2\, d\Omega.
\tag{7.175}
\end{equation*}
这里，积分取在空间无穷远处的二维球面 $S^2_\infty$ 上，而 $n^\mu$ 是垂直于 $S^2_\infty$ 的单位类空矢量。在极坐标 $\{t, r, \theta, \phi\}$ 中，法矢量的分量为
\begin{equation*}
n^\mu = (0,\, 1,\, 0,\, 0).
\tag{7.176}
\end{equation*}

我们想用 $t_{\mu\nu}$ 的表达式 (7.164) 来计算功率 $P$。面临的第一个问题是，这个表达式是用横无迹微扰写下的，而四极公式 (7.140) 却是用 Lorenz 规范下无迹反转微扰的空间分量 $h_{ij}$ 写下的。最简单的步骤（虽然它也没那么简单）是先把 $h_{ij}$ 转换到横无迹规范——由于我们关心的是波在远离其发射源的真空中的行为，这样做是允许的——代入 $t_{\mu\nu}$ 的公式，然后再转换回非横无迹的形式。让我们看看这一套流程是如何进行的。

我们首先引入（空间的）投影张量
\begin{equation*}
P_{ij} = \delta_{ij} - n_i n_j,
\tag{7.177}
\end{equation*}
它把张量分量投影到与单位矢量 $n^i$ 正交的曲面上。（更详细的讨论见附录 D。）在我们的情形中，我们取 $n^i$ 沿波的传播方向，这样 $P_{ij}$ 便投影到空间无穷远处的二维球面上。我们可以用这个投影张量，按下式构造对称空间张量 $X_{ij}$ 的横无迹版本：
\begin{equation*}
X^{\mathrm{TT}}_{ij} = \left(P_i{}^k P_j{}^l - \frac{1}{2}P_{ij}P^{kl}\right)X_{kl}.
\tag{7.178}
\end{equation*}
你可以自行验证 $X^{\mathrm{TT}}_{ij}$ 确实是横向且无迹的。由于它无迹，$\bar{h}^{\mathrm{TT}}_{ij}$ 就等于原来的微扰 $h^{\mathrm{TT}}_{ij}$；代入四极公式 (7.140)，我们得到
\begin{equation*}
h^{\mathrm{TT}}_{ij} = \bar{h}^{\mathrm{TT}}_{ij} = \frac{2G}{r}\frac{d^2 I^{\mathrm{TT}}_{ij}}{dt^2}(t - r),
\tag{7.179}
\end{equation*}
其中四极矩的横无迹部分同样按 (7.178) 构造。事实上，由 (7.138) 定义的四极矩并不是用来表达所生成的波的最方便的量，因为它涉及一个
% ---- 书页 313 / p328 ----
对能量密度的积分，而这个积分可能难以确定。我们可以改用\textbf{约化四极矩}（reduced quadrupole moment）
\begin{equation*}
J_{ij} = I_{ij} - \frac{1}{3}\delta_{ij}\delta^{kl}I_{kl},
\tag{7.180}
\end{equation*}
它正是 $I_{ij}$ 的无迹部分。约化四极矩有一个很好的性质：它是 Newton 势多极展开中 $r^{-3}$ 项的系数，
\begin{equation*}
\Phi = -\frac{GM}{r} - \frac{G}{r^3}D_i x^i - \frac{3G}{2r^5}J_{ij}x^i x^j + \cdots,
\tag{7.181}
\end{equation*}
因此对于现实的源更容易近似。（这里 $D_i$ 是偶极矩，$D_i = \int T^{00}x^i\, d^3x$。）当然，四极矩的横无迹部分与约化（也就是无迹）四极矩的横无迹部分是一样的，所以 (7.179) 变为
\begin{equation*}
h^{\mathrm{TT}}_{ij} = \frac{2G}{r}\frac{d^2 J^{\mathrm{TT}}_{ij}}{dt^2}(t - r).
\tag{7.182}
\end{equation*}

为了计算功率，我们关心的是 $t_{0\mu}n^\mu = t_{0r}$。由于四极矩只依赖于推迟时间（retarded time）$t_r = t - r$，我们有
\begin{align*}
\partial_0 h^{\mathrm{TT}}_{ij} &= \frac{2G}{r}\frac{d^3 J^{\mathrm{TT}}_{ij}}{dt^3},\\
\partial_r h^{\mathrm{TT}}_{ij} &= -\frac{2G}{r}\frac{d^3 J^{\mathrm{TT}}_{ij}}{dt^3} - \frac{2G}{r^2}\frac{d^2 J^{\mathrm{TT}}_{ij}}{dt^2}\\
&\approx -\frac{2G}{r}\frac{d^3 J^{\mathrm{TT}}_{ij}}{dt^3},
\tag{7.183}
\end{align*}
其中我们已经去掉了 $r^{-2}$ 项，因为我们关心的是 $r \to \infty$ 的极限。于是，能量-动量张量的重要分量为
\begin{equation*}
t_{0r} = -\frac{G}{8\pi r^2}\left\langle\left(\frac{d^3 J^{\mathrm{TT}}_{ij}}{dt^3}\right)\left(\frac{d^3 J_{\mathrm{TT}}^{ij}}{dt^3}\right)\right\rangle.
\tag{7.184}
\end{equation*}
下一步是从横无迹部分转换回 $J_{ij}$。应用 (7.178) 以及一些繁琐的代数运算，可以直接证明
\begin{equation*}
X^{\mathrm{TT}}_{ij}X_{\mathrm{TT}}^{ij} = X_{ij}X^{ij} - 2X_i{}^j X^{ik}n_j n_k + \frac{1}{2}X^{ij}X^{kl}n_i n_j n_k n_l - \frac{1}{2}X^2 + XX^{ij}n_i n_j,
\tag{7.185}
\end{equation*}
其中 $X = \delta^{ij}X_{ij}$。由于 $J_{ij}$ 是无迹的，我们有
\begin{equation*}
J^{\mathrm{TT}}_{ij}J_{\mathrm{TT}}^{ij} = J_{ij}J^{ij} - 2J_i{}^j J^{ik}n_j n_k + \frac{1}{2}J^{ij}J^{kl}n_i n_j n_k n_l,
\tag{7.186}
\end{equation*}

% ---- 书页 314 / p329 ----

而功率为
\begin{align*}
P = -\frac{G}{8\pi}\int_{S^2_\infty}\bigg\langle \frac{d^3 J_{ij}}{dt^3}\frac{d^3 J^{ij}}{dt^3} &- 2\frac{d^3 J_i{}^j}{dt^3}\frac{d^3 J^{ik}}{dt^3}n_j n_k\\
&+ \frac{1}{2}\frac{d^3 J^{ij}}{dt^3}\frac{d^3 J^{kl}}{dt^3}n_i n_j n_k n_l\bigg\rangle\, d\Omega.
\tag{7.187}
\end{align*}
要求出这个表达式的值，最好换回到空间的笛卡尔坐标，其中 $n^i = x^i/r$。四极张量与角坐标无关，因为它们是由对全空间的积分定义的。于是我们可以把它们移到积分号外，并利用恒等式
\begin{align*}
\int d\Omega &= 4\pi\\
\int n_i n_j\, d\Omega &= \frac{4\pi}{3}\delta_{ij}\\
\int n_i n_j n_k n_l\, d\Omega &= \frac{4\pi}{15}\left(\delta_{ij}\delta_{kl} + \delta_{ik}\delta_{jl} + \delta_{il}\delta_{jk}\right).
\tag{7.188}
\end{align*}
当一切尘埃落定，功率的表达式便塌缩为
\begin{equation*}
P = -\frac{G}{5}\left\langle \frac{d^3 J_{ij}}{dt^3}\frac{d^3 J^{ij}}{dt^3}\right\rangle,
\tag{7.189}
\end{equation*}
这里我们应当记住，四极矩是在推迟时间 $t_r = t - r$ 处取值的。我们的公式带有一个负号，因为它表示的是能量变化的速率，而辐射源将不断损失能量。

对于由 (7.148) 所表示的双星系统，约化四极矩为
\begin{equation*}
J_{ij} = \frac{MR^2}{3}
\begin{pmatrix}
(1+3\cos 2\Omega t) & 3\sin 2\Omega t & 0\\
3\sin 2\Omega t & (1-3\cos 2\Omega t) & 0\\
0 & 0 & -2
\end{pmatrix},
\tag{7.190}
\end{equation*}
它对时间的三阶导数于是为
\begin{equation*}
\frac{d^3 J_{ij}}{dt^3} = 8MR^2\Omega^3
\begin{pmatrix}
\sin 2\Omega t & -\cos 2\Omega t & 0\\
-\cos 2\Omega t & -\sin 2\Omega t & 0\\
0 & 0 & 0
\end{pmatrix}.
\tag{7.191}
\end{equation*}
于是双星辐射的功率为
\begin{equation*}
P = -\frac{128}{5}GM^2R^4\Omega^6,
\tag{7.192}
\end{equation*}
或者，利用 (7.144) 的频率表达式，
\begin{equation*}
P = -\frac{2}{5}\frac{G^4M^5}{R^5}.
\tag{7.193}
\end{equation*}

% ---- 书页 315 / p330 ----

当然，通过发射引力辐射而损失能量的现象已经被观测到了。1974 年，Hulse 和 Taylor 发现了一个双星系统 PSR~1913+16，其中的两颗星都非常小，因此经典效应可以忽略，或者至少是可控的，而且其中一颗是脉冲星。轨道周期为八小时，按天体物理的标准衡量极其短。其中一颗星是脉冲星这个事实提供了一个非常精确的时钟，借助它便可以测量系统损失能量时周期的变化。这个结果与广义相对论关于引力辐射导致能量损失的预言相一致。

\section{引力波的探测}

当代引力物理和天体物理学最优先的目标之一，是直接探测引力辐射。（所谓“直接”，我们的意思是“通过观测引力波对检验物体的影响”，与之相对的是观测能量损失的间接效应，就像脉冲双星中的情形。）完全有理由相信，这样的探测很快就会实现，或者是在已经建成的引力波天文台中，或者是在近期正在规划的引力波天文台中。当然，一旦我们探测到了引力辐射，目标就会立刻变成从观测中提取有用的天体物理信息。我们目前对太阳系之外宇宙的认识，几乎完全来自对电磁辐射的观测，外加一些来自中微子和宇宙线的输入；引力波天体物理学的到来，将打开一扇通往遥远宇宙中高能现象的全新窗口。\footnote{关于引力波天体物理学的概览，见 S.A.~Hughes, S.~Márka, P.L.~Bender 与 C.J.~Hogan, “New physics and astronomy with the new gravitational-wave observatories,” \texttt{http://arxiv.org/astro-ph/0110349}。}

在讨论我们可以如何着手探测天体物理引力波之前，我们应该想一想哪些源最容易被观测到。第一个重要的认识是：产生可观引力辐射所需的条件，与产生电磁辐射所需的条件非常不同。这个差别可以追溯到这样一个事实：引力波由大质量物体的整体运动（bulk motion）产生，而电磁波（通常）由单个粒子的不相干激发产生。因此，电磁辐射可以由一个整体上静态的源产生，比如一颗恒星，这对天文学家来说是一大优势。然而，引力波是由大块运动着的物质相干地产生的（物质中的每个粒子都以相同的方式对波有所贡献），这可以部分弥补静态源无法发射引力波的缺憾。

因此，我们需要具有可观整体运动的大质量源。作为一个简单的例子，考虑 7.5 节中的双星系统，其中两颗星的质量都是 $M$，轨道半径为 $R$。我们将稍微作一点弊，把 Newton 的轨道参数公式用到广义相对论已经开始变得重要的区域，不过对于数量级的估计这已经足够了。相关参数可以归结为 Schwarzschild 半径 $R_S = 2GM/c^2$、
% ---- 书页 316 / p331 ----
轨道半径 $R$，以及我们与双星之间的距离 $r$。（现在我们将把因子 $c$ 显式地恢复，以便与实验比较。）用这些量来表示，轨道的频率——也即所产生的引力波的频率——近似为
\begin{equation*}
f = \frac{\Omega}{2\pi} \sim \frac{cR_S^{1/2}}{10R^{3/2}}.
\tag{7.194}
\end{equation*}
利用 (7.149) 中关于所产生的微扰的公式，我们可以估计接收到的引力波振幅为
\begin{equation*}
h \sim \frac{R_S^2}{rR}.
\tag{7.195}
\end{equation*}
让我们看看，这对我们有望观测到的那一类源意味着什么。一个范例性的例子是黑洞/黑洞双星的并合（coalescence）。对于典型的参数，我们可以取两个黑洞都是 10 倍太阳质量，双星处于约 100 Mpc 的宇宙学距离上，而两颗子星相距为其 Schwarzschild 半径的十倍：
\begin{gather*}
R_S \sim 10^6\,\mathrm{cm}\\
R \sim 10^7\,\mathrm{cm}\\
r \sim 10^{26}\,\mathrm{cm}.
\tag{7.196}
\end{gather*}
这样的源因而具有如下特征量级：
\begin{equation*}
f \sim 10^2\,\mathrm{s}^{-1}, \qquad h \sim 10^{-21}.
\tag{7.197}
\end{equation*}
如果我们还想有希望探测到具有这些参数的双星并合，我们就需要对 100 Hz 附近的频率以及 $10^{-21}$ 或更小量级的应变足够敏感。

所幸的是，这些参数处于我们实验能力（在许多科学家英勇的努力之下）所能及的范围之内。目前正在考虑的引力波探测方案中最有前途的技术是干涉测量法，因此这里我们将只讨论干涉仪，尽管完全可以设想，将来会发明出具有更好灵敏度的新技术。

回想一下，引力波经过时的物理效应是轻微地扰动自由下落的质量之间的相对位置。如果两个检验质量（test masses）相距 $L$，它们之间距离的变化大致为
\begin{equation*}
\frac{\delta L}{L} \sim h.
\tag{7.198}
\end{equation*}
设想我们打算建造一座检验物体相距公里量级的天文台。那么，要探测振幅为
% ---- 书页 317 / p332 ----
$h\sim 10^{-21}$ 量级的波，就需要对量级为
\begin{equation*}
\delta L \sim 10^{-16}\left(\frac{h}{10^{-21}}\right)\left(\frac{L}{\mathrm{km}}\right)\,\mathrm{cm}
\tag{7.199}
\end{equation*}
的变化具有足够的灵敏度。把它与一个典型原子的大小比较一下——后者由玻尔半径（Bohr radius）给出，
\begin{equation*}
a_0 \sim 5\times 10^{-9}\,\mathrm{cm},
\tag{7.200}
\end{equation*}
或者，也不妨与一个典型原子核的大小比较——大约一个费米（Fermi），
\begin{equation*}
1\,\mathrm{fm} = 10^{-13}\,\mathrm{cm}.
\tag{7.201}
\end{equation*}
我们在这里反复强调的要点是：一座切实可行的地面引力波天文台，必须对距离变化足够敏感——要比任何可以设想的检验质量所必须由之构成的那些原子的大小还要小得多。

%FIG{7.11}: {引力波干涉仪的示意图。}

激光干涉仪提供了一条克服测量如此微小扰动之困难的途径。考虑图 7.11 所描绘的示意图装置。一束激光（典型特征波长 $\lambda \sim 10^{-4}\,\mathrm{cm}$）被导向一个分束器（beamsplitter），分束器把光子送进两条抽成真空的、长度为 $L$ 的管道。腔的末端是检验质量，用悬挂在摆上的反射镜来表示。光实际上是在分束器附近的部分反射镜（partially-reflective mirrors）之间来回反射的，所以一个典型的光子在回到分束器并被导入
% ---- 书页 318 / p333 ----
光电二极管之前，会在腔内上下往返约 100 次。整个系统的安排是这样的：如果检验质量完全静止，返回的两束光将发生相消干涉，光电二极管接收不到任何信号。正如我们已经看到的，引力波经过的效应是把相互正交的长度朝相反的方向扰动，这会引起激光脉冲的相移，从而破坏相消干涉。在腔臂中往返 100 次的过程中，累积的相移将为
\begin{equation*}
\delta\phi \sim 200\left(\frac{2\pi}{\lambda}\right)\delta L \sim 10^{-9},
\tag{7.202}
\end{equation*}
其中用 200 而不是 100，反映的是两臂中的相移会相加这一事实。这样微小的相移是可以测量的，只要光子数 $N$ 足够大，足以克服“散粒噪声”（shot noise）；具体地说，只要 $\sqrt{N} > \delta\phi$。

建造足够“安静”且灵敏的引力波天文台所面临的技术挑战，正在多个不同的地方被攻克，包括美国（LIGO）、意大利（Virgo）、德国（GEO）、日本（TAMA）和澳大利亚（ACIGA）。LIGO（激光干涉引力波天文台，Laser Interferometric Gravitational-Wave Observatory）是目前最先进的探测器；它由两个设施组成（一个在华盛顿州，一个在路易斯安那州），每个都有四公里长的臂。单一的引力波天文台将无法对源在天空中的位置定位；要完成这一任务，多个探测器至关重要（这对于验证一个表观信号是否真实同样关键）。

基本的噪声来源限制了地面天文台探测低频引力波的能力。图 7.12 展示了两种截然不同的设计——LIGO 这样的地面天文台，以及 LISA（激光干涉空间天线，Laser Interferometer Space Antenna）这样的空间任务——的灵敏度区域随频率的变化。LISA 背后的一般原理与其他任何干涉仪相同，但其实施方案则（如果它真的被建造的话，或者说将会）截然不同。目前的设计设想三个航天器绕太阳运行，位于地球后方约三千万公里处，彼此相距五百万公里。由于间隔大得多，LISA 对 $10^{-2}\,\mathrm{Hz}$ 附近的频率敏感。这张图中所描绘的灵敏度只能视为示意性的，因为它们依赖于积分时间和其他因素。

许多潜在的噪声源摆在了引力波天文学家的面前。对于地面天文台，低频处占主导的效应通常是地震噪声（seismic noise），高频处的效应来自光子散粒噪声，中间频率处则来自热噪声。地面探测器的改进版本也许能够补偿低频的地震噪声，但会遇到来自大气现象或附近经过的物体（比如汽车）的引力梯度所带来的不可约噪声。卫星天文台当然对这类效应免疫。相反，其根本限制预期来自测量航天器之间距离变化时的误差（更确切地说，是航天器内部受到屏蔽的检验质量之间的距离变化），以及航天器的非引力加速度。

% ---- 书页 319 / p334 ----

%FIG{7.12}: {代表性的地面（LIGO）与空间（LISA）引力波天文台的灵敏度随频率的变化，以及各种可能的源所预期的信号。图取自 LISA 合作组主页（\texttt{http://lisa.jpl.nasa.gov/}）。}

最后，我们可以对引力波天文台各种可能的源作一个十分简要的概述。我们已经提到过各种类型的致密双星。对于地面天文台，这类源要等到非常接近并合时才会变得可见，而且前提是两颗子星的质量足够大（中子星或黑洞）。根据我们对这类系统的认识作外推，在几百 Mpc 的距离之内每年也许会有几次并合。另一个有希望的可能性是大质量恒星的核心坍缩（core collapse），即产生超新星的过程。虽然完全球对称的坍缩不会产生任何引力波，但现实的事件预期会受到破坏这种对称性的不稳定性的影响。一个令人兴奋的前景是用普通望远镜和引力波天文台对超新星进行协同观测。最后，在地面天文台可能的源之中，还有周期性的源，比如（非完全轴对称的）旋转中子星。这类源的振幅预期较小，但未必完全超出先进探测器的可及范围。

对空间探测器而言，有趣的源则有些不同。最重要的是，我们银河系中已知的双星群体必将提供强度可探测的引力波信号。事实上，无法分辨的双星会构成探测器的一种混淆噪声（confusion noise）的来源，因为我们无法从背景中挑出个别的低强度源。尽管如此，众多较高强度的源应当是容易观测到的。此外，超大质量黑洞（大于 $1000\,M_\odot$，比如位于星系中心的那些）演化过程中的各种
% ---- 书页 320 / p335 ----
过程也会带来有趣的源：这类天体的形成、它们通过吸积更小天体而实现的后续成长，以及多个超大质量黑洞可能的并合。追踪一个太阳质量的黑洞绕行并最终落入超大质量黑洞时其引力波信号的演化，将使我们能够精确绘制时空度规，从而为广义相对论提供一种新颖的检验。

除了局域源产生的波之外，我们还面临着随机引力波背景（stochastic gravitational-wave backgrounds）的可能性。我们指的是一个各向同性的引力波集合，它们或许产生于早期宇宙，其特征是具有随频率平滑变化的功率谱。一种可能性是由暴胀产生的近乎无标度（scale-free）的引力波谱，这将在第 8 章讨论。这样的波几乎不可能在地面上被直接探测到（也许比先进探测器的探测能力还要低大约五个数量级），甚至连 LISA 也无能为力，但可以设想它们会被下一代空间任务观测到。更可能的情况是，任何这样的波将首先在宇宙微波背景的偏振中显现。然而，另一种可能性是从剧烈的（一级）相变中产生原初引力波。这样的波将具有一个带有明确峰值的频谱，峰频与相变温度 $T$ 的关系为
\begin{equation*}
f_{\mathrm{peak}} \sim 10^{-3}\left(\frac{T}{1000\,\mathrm{GeV}}\right)\,\mathrm{Hz}.
\tag{7.203}
\end{equation*}
这样，一级电弱相变（$T \sim 200\,\mathrm{GeV}$）恰好落在 LISA 潜在可观测的频带之内。这一点尤其耐人寻味，因为某些重子产生（baryogenesis）模型要求在这个尺度上存在一个强相变；想一想我们居然能通过一个引力实验学到一些关于电弱物理的重要东西，这是相当诱人的。

\section{习题}

\textbf{1.}\quad 证明：拉格朗日量 (7.9) 给出 Einstein 方程的线性化版本。

\textbf{2.}\quad 考虑一个薄球形物质壳，质量为 $M$、半径为 $R$，以角速度 $\Omega$ 缓慢旋转。

\textbf{(a)}\quad 证明引力电场 $\vec{G}$ 为零，并用 $M$、$R$ 和 $\Omega$ 计算引力磁场 $\vec{H}$。

\textbf{(b)}\quad 壳所产生的非零引力磁场会导致惯性系被拖曳，这就是\textbf{Lense--Thirring 效应}（Lense--Thirring effect）。计算位于壳中心的一个自由下落观者的转动（相对于由背景 Minkowski 度规定义的惯性系）。换言之，计算位于中心处的一个被平行移动的矢量的空间分量的进动。

% ---- 书页 321 / p336 ----

\textbf{3.}\quad Fermat 原理指出，光线沿耗时最少的路径传播。对于折射率为 $n(\mathbf{x})$ 的介质，这等价于对时间
\begin{equation*}
t = \int n(\mathbf{x})\left[\delta_{ij}\, dx^i dx^j\right]^{1/2}
\tag{7.204}
\end{equation*}
沿路径求极值。证明：当折射率取 $n = 1 - 2\Phi$ 时，Fermat 原理给出受 Newton 势微扰的时空中光子的正确运动方程。

\textbf{4.}\quad 证明：Lorenz 规范条件 $\partial_\mu\bar{h}^{\mu\nu} = 0$ 等价于\textbf{谐和规范}（harmonic gauge）条件。这个规范由
\begin{equation*}
\Box x^\mu = 0,
\tag{7.205}
\end{equation*}
定义，其中每个坐标 $x^\mu$ 都被当作时空上的标量函数。（任何满足 $\Box f = 0$ 的函数都称为“谐和函数”（harmonic function）。）

\textbf{5.}\quad 在第 3 章的习题中，我们引入了度规
\begin{equation*}
ds^2 = -(\mathrm{d}u\,\mathrm{d}v + \mathrm{d}v\,\mathrm{d}u) + a^2(u)\mathrm{d}x^2 + b^2(u)\mathrm{d}y^2,
\tag{7.206}
\end{equation*}
其中 $a$ 和 $b$ 是 $u$ 的未具体指定的函数。对于适当的函数 $a$ 和 $b$，它表示一个\emph{精确}的引力平面波。

\textbf{(a)}\quad 计算这个度规的 Christoffel 符号和 Riemann 张量。

\textbf{(b)}\quad 利用真空中的 Einstein 方程，推导 $a(u)$ 和 $b(u)$ 所满足的方程。

\textbf{(c)}\quad 证明可以找到一个精确解，其中 $a$ 和 $b$ 都由一个\emph{任意}函数 $f(u)$ 确定。

\textbf{6.}\quad 两个质量为 $M$ 的物体在事件 $(0,0,0,0)$ 处迎面对撞。在遥远的过去 $t \to -\infty$，两物体从 $x \to \pm\infty$ 处由静止出发。

\textbf{(a)}\quad 利用 Newton 理论，证明 $x(t) = \pm(9Mt^2/8)^{1/3}$。

\textbf{(b)}\quad 对多大的间距，Newton 近似是合理的？

\textbf{(c)}\quad 计算 $(x, y, z) = (0, R, 0)$ 处的 $h^{\mathrm{TT}}_{xx}(t)$。

\textbf{7.}\quad 引力波可以通过监测两个自由飞行的质量之间的距离来探测。如果其中一个质量配有激光和精确的时钟，另一个配有好的反射镜，那么通过对一束激光脉冲往返一趟所需时间的计时，就可以测量两个质量之间的距离。要使你的探测器对形如
\begin{equation*}
ds^2 = -\mathrm{d}t^2 + \left[1 + A\cos(\omega(t-z))\right]\mathrm{d}x^2 + \left[1 - A\cos(\omega(t-z))\right]\mathrm{d}y^2 + \mathrm{d}z^2?
\end{equation*}
的平面波产生最大的响应，你会希望探测器如何取向？如果两个质量的平均间距为 $L$，波造成的脉冲到达时间的最大变化是多少？哪些频率 $\omega$ 会探测不到？

\textbf{8.}\quad 当两个质量彼此散射时，会产生轫致辐射（\emph{bremsstrahlung}）的引力类比。考虑一个质量为 $m$ 的小质量以碰撞参数 $b$、总能量 $E = 0$ 散射到一个大质量 $M$ 上时会发生什么。取 $M \gg m$ 和 $M/b \ll 1$。小质量的运动可以用 Newton 物理来描述，因为 $M/b \ll 1$。如果轨道位于 $(x, y)$ 平面内，且大质量位于

% ---- 书页 322 / p337 ----

$(x, y, z) = (0, 0, 0)$ 处，试计算 $(x, y, z) = (0, 0, r)$ 处两个偏振的引力波振幅。由于运动不是周期性的，引力波将是爆发式的，并由许多不同的频率组成。根据物理上的理由，你预期主频是多少？估计系统辐射的总能量。它与这个小质量的峰值动能相比如何？

\emph{提示}：轨道的解可以在 Goldstein (2002) 中找到。解为
\begin{gather*}
r = \frac{2b}{1+\cos\theta},\\
t = \sqrt{\frac{2b^3}{M}}\left(\tan\frac{\theta}{2} + \frac{1}{3}\tan^3\frac{\theta}{2}\right).
\end{gather*}
时间的范围是 $t = (-\infty, \infty)$。与其使用上述 $\theta(t)$ 的隐式解，你或许更愿意用
\begin{equation*}
\dot{\theta} = \sqrt{\frac{M}{8b^3}}\,(1+\cos\theta)^2.
\end{equation*}

\textbf{9.}\quad 验证：引力波能量-动量张量的表达式 (7.165) 在规范变换 $h_{\mu\nu} \to h_{\mu\nu} + 2\partial_{(\mu}\xi_{\nu)}$ 下不变。

\textbf{10.}\quad 证明：引力微扰总能量的积分表达式 (7.173) 与空间超曲面 $\Sigma$ 无关。
